
\documentclass[11pt]{article}

\usepackage[margin=1in]{geometry}
\usepackage{amsmath, amssymb, amsthm, mathtools}
\usepackage{enumitem}
\usepackage{xcolor}
\usepackage{hyperref}
\usepackage{microtype}

\hypersetup{
    colorlinks=true,
    linkcolor=blue!60!black,
    urlcolor=blue!60!black,
    citecolor=blue!60!black
}

\newtheorem{theorem}{Theorem}[section]
\newtheorem{lemma}[theorem]{Lemma}
\newtheorem{proposition}[theorem]{Proposition}
\newtheorem{corollary}[theorem]{Corollary}
\newtheorem{definition}[theorem]{Definition}
\newtheorem{remark}[theorem]{Remark}

\newcommand{\N}{\mathbb{N}}
\newcommand{\R}{\mathbb{R}}
\newcommand{\one}{\mathbf{1}}
\newcommand{\State}{\mathcal{X}}
\newcommand{\Success}{\mathcal{S}}
\newcommand{\Failure}{\mathcal{F}}
\newcommand{\Active}{\mathcal{A}}
\newcommand{\Actions}{\mathcal{U}}
\newcommand{\Prob}{\mathbb{P}}
\newcommand{\E}{\mathbb{E}}
\newcommand{\rank}{\rho}
\newcommand{\weight}{\omega}
\newcommand{\impl}{\mathrm{impl}}
\newcommand{\opt}{\mathrm{opt}}

\title{
    \vspace{-1.5em}
    Optimal Horse-Combination Probability\\
    \large A Bellman Correctness Proof for a Recursive Tier-Upgrade Calculator
}
\author{Prepared proof draft}
\date{\today}

\begin{document}

\maketitle

\begin{abstract}
This document formalizes the recursive horse-combination probability calculator as a finite Markov decision process and proves the correctness of its core value recursion.  A state records the number of horses in tiers \(1,2,3,4\).  The objective is to maximize the probability of eventually reaching a state containing at least one tier-4 horse.  The main result shows that the recursive function
\[
    P(s)=\max_{a\in \Actions(s)} \sum_{s'} K(s' \mid s,a)P(s')
\]
with terminal boundary values \(P(s)=1\) on success states and \(P(s)=0\) on terminal failure states computes the optimal success probability.

The proof uses a well-founded rank,
\[
    \rank(s)=t_1+t_2+t_3+t_4,
\]
which strictly decreases by one after every legal nonterminal combination.  This makes the state-transition graph acyclic and allows a strong-induction proof of Bellman correctness.  We also prove the minimal-resource lemma: fewer than eight tier-1 horses cannot reach tier 4, and exactly eight tier-1 horses succeed with probability \((1/5)^7\).  Finally, we record a careful interpretation of empirical finite-window scaling: a fitted law such as \(P(X)\propto X^\alpha\) may describe a finite range, but no positive power law can hold globally because probabilities are bounded by \(1\).
\end{abstract}

\tableofcontents

\section{Purpose}

The original implementation computes the probability of obtaining at least one tier-4 horse by recursively evaluating all legal horse-combination actions and selecting the action with maximal expected success probability.  This document supplies the mathematical proof spine behind that implementation.

There are three layers:

\begin{enumerate}[label=\textbf{Layer \arabic*.}]
    \item \textbf{Formal model.}  Define states, legal actions, transition probabilities, terminal success, and terminal failure.
    \item \textbf{Correctness theorem.}  Prove that the recursive value function equals the optimal probability of eventual success.
    \item \textbf{Auxiliary proof notes.}  Prove the eight-horse base case and distinguish exact theorem claims from empirical finite-range scaling observations.
\end{enumerate}

The central point is that the recursion is not justified merely by memoization.  Memoization is an implementation optimization.  The mathematical reason the recursion is well-founded is that every legal combination reduces the total horse count by exactly one.

\section{State Space}

\begin{definition}[Horse state]
A horse state is a vector
\[
    s=(t_1,t_2,t_3,t_4)\in \N^4,
\]
where \(t_i\) is the number of horses in tier \(i\).  Tier \(4\) is the goal tier.
\end{definition}

Let
\[
    e_1=(1,0,0,0),\quad
    e_2=(0,1,0,0),\quad
    e_3=(0,0,1,0),\quad
    e_4=(0,0,0,1)
\]
denote the standard basis vectors.

\begin{definition}[Success states]
The success set is
\[
    \Success=\{s=(t_1,t_2,t_3,t_4)\in \N^4:t_4\ge 1\}.
\]
When \(s\in\Success\), the process has already achieved the objective.
\end{definition}

\begin{definition}[Legal action]
For \(i\in\{1,2,3\}\), action \(a_i\) means ``combine two tier-\(i\) horses.''  It is legal at state \(s\) exactly when
\[
    t_i\ge 2.
\]
The legal action set is
\[
    \Actions(s)=\{a_i:i\in\{1,2,3\},\ t_i\ge 2\}.
\]
\end{definition}

\begin{definition}[Failure states]
The terminal failure set is
\[
    \Failure=\{s\in \N^4:t_4=0,\ \Actions(s)=\varnothing\}.
\]
Equivalently,
\[
    \Failure=\{s=(t_1,t_2,t_3,0):t_1,t_2,t_3\le 1\}.
\]
\end{definition}

\begin{definition}[Active states]
The active set is
\[
    \Active=\N^4\setminus(\Success\cup\Failure).
\]
These are precisely the states for which the process has not succeeded and at least one legal combination remains.
\end{definition}

\section{Transition Rules}

The transition probabilities are
\[
    p_+ = \frac15,\qquad
    p_0 = \frac12,\qquad
    p_- = \frac3{10}.
\]
Here \(p_+\) is the probability of upgrading to the next tier, \(p_0\) is the probability of remaining at the same tier, and \(p_-\) is the probability of downgrading by one tier.  For tier \(1\), downgrade is absorbed into the tier-\(1\) outcome because there is no tier below tier \(1\).

\subsection{Combining two tier-1 horses}

If \(a_1\) is legal at \(s\), then
\[
    s \longrightarrow s-2e_1+e_2
    \quad\text{with probability}\quad \frac15,
\]
and
\[
    s \longrightarrow s-e_1
    \quad\text{with probability}\quad \frac45.
\]
The second transition represents the combined same-tier/downgrade case: two tier-1 horses are consumed and one tier-1 horse remains.

\subsection{Combining two tier-2 horses}

If \(a_2\) is legal at \(s\), then
\[
    s \longrightarrow s-2e_2+e_3
    \quad\text{with probability}\quad \frac15,
\]
\[
    s \longrightarrow s-e_2
    \quad\text{with probability}\quad \frac12,
\]
and
\[
    s \longrightarrow s+e_1-2e_2
    \quad\text{with probability}\quad \frac3{10}.
\]

\subsection{Combining two tier-3 horses}

If \(a_3\) is legal at \(s\), then
\[
    s \longrightarrow s-2e_3+e_4
    \quad\text{with probability}\quad \frac15,
\]
\[
    s \longrightarrow s-e_3
    \quad\text{with probability}\quad \frac12,
\]
and
\[
    s \longrightarrow s+e_2-2e_3
    \quad\text{with probability}\quad \frac3{10}.
\]

The first tier-3 outcome is a success transition, since it creates a tier-4 horse.

\section{Policies and Optimal Value}

\begin{definition}[Policy]
A policy \(\pi\) is a rule that assigns a legal action to every active state encountered by the process.  Since each legal move decreases total horse count, the process has finite horizon from every finite initial state.
\end{definition}

\begin{definition}[Optimal success probability]
For any state \(s\), let
\[
    V^\opt(s)
\]
denote the supremum, over all admissible policies \(\pi\), of the probability of eventually reaching \(\Success\) before reaching \(\Failure\), starting from \(s\).
\end{definition}

Since the action set is finite and every play terminates in finitely many moves, the supremum is attained.  Thus \(V^\opt\) is a maximum value rather than only a supremal value.

\section{Well-Founded Rank}

\begin{definition}[Total-horse rank]
Define
\[
    \rank(s)=t_1+t_2+t_3+t_4.
\]
\end{definition}

\begin{lemma}[Every legal move decreases rank]
Let \(s\in\Active\), and let \(a\in\Actions(s)\).  Every possible successor state \(s'\) produced by action \(a\) satisfies
\[
    \rank(s')=\rank(s)-1.
\]
\end{lemma}

\begin{proof}
Every legal action combines two horses and returns exactly one horse.  Thus the total number of horses decreases by one.  Checking the transition formulas explicitly gives the same result.

For \(a_1\),
\[
    \rank(s-2e_1+e_2)=\rank(s)-1
\]
and
\[
    \rank(s-e_1)=\rank(s)-1.
\]

For \(a_2\),
\[
    \rank(s-2e_2+e_3)=\rank(s)-1,
\]
\[
    \rank(s-e_2)=\rank(s)-1,
\]
and
\[
    \rank(s+e_1-2e_2)=\rank(s)-1.
\]

For \(a_3\),
\[
    \rank(s-2e_3+e_4)=\rank(s)-1,
\]
\[
    \rank(s-e_3)=\rank(s)-1,
\]
and
\[
    \rank(s+e_2-2e_3)=\rank(s)-1.
\]

Therefore every legal transition decreases rank by exactly one.
\end{proof}

\begin{corollary}[No transition cycles]
The directed state-transition graph contains no directed cycle among nonterminal states.
\end{corollary}

\begin{proof}
Along every nonterminal transition, \(\rank\) strictly decreases.  A directed cycle would require returning to the original rank, which is impossible under strict decrease.
\end{proof}

\begin{corollary}[Finite termination]
Starting from any finite state \(s\), every trajectory reaches either \(\Success\) or \(\Failure\) after at most \(\rank(s)\) legal combinations.
\end{corollary}

\begin{proof}
The rank is a nonnegative integer and decreases by one after every legal nonterminal move.  Therefore the process cannot continue for more than \(\rank(s)\) legal moves.  It must either reach a success state earlier or eventually have no legal action remaining, which is a failure state if no tier-4 horse has been produced.
\end{proof}

\section{Bellman Recurrence}

\begin{definition}[Bellman operator]
For any function \(W:\N^4\to[0,1]\), define the action-value operator
\[
    Q_W(s,a)=\sum_{s'}K(s'\mid s,a)W(s'),
\]
where \(K(s'\mid s,a)\) is the transition kernel described above.  For active states, define
\[
    (TW)(s)=\max_{a\in\Actions(s)}Q_W(s,a).
\]
\end{definition}

\begin{lemma}[Bellman optimality equation]
The optimal value function satisfies
\[
    V^\opt(s)=1 \quad\text{for }s\in\Success,
\]
\[
    V^\opt(s)=0 \quad\text{for }s\in\Failure,
\]
and
\[
    V^\opt(s)
    =
    \max_{a\in\Actions(s)}
    \sum_{s'}K(s'\mid s,a)V^\opt(s')
    \quad\text{for }s\in\Active.
\]
\end{lemma}

\begin{proof}
If \(s\in\Success\), the objective has already been achieved, so the success probability is \(1\).  If \(s\in\Failure\), no tier-4 horse exists and no legal move remains, so the success probability is \(0\).

Now let \(s\in\Active\).  Any admissible policy must choose some first action \(a\in\Actions(s)\).  Conditional on the successor state \(s'\), the remaining success probability cannot exceed \(V^\opt(s')\).  Therefore the success probability of any policy whose first action is \(a\) is at most
\[
    \sum_{s'}K(s'\mid s,a)V^\opt(s').
\]
Since the first action must be legal, every policy has success probability at most
\[
    \max_{a\in\Actions(s)}
    \sum_{s'}K(s'\mid s,a)V^\opt(s').
\]

Conversely, choose an action \(a^*\in\Actions(s)\) attaining the maximum.  Such an action exists because \(\Actions(s)\) is finite.  For each successor state \(s'\), follow an optimal policy from \(s'\).  This patched policy achieves
\[
    \sum_{s'}K(s'\mid s,a^*)V^\opt(s')
    =
    \max_{a\in\Actions(s)}
    \sum_{s'}K(s'\mid s,a)V^\opt(s').
\]
Thus equality holds.
\end{proof}

\section{Recursive Implementation as a Mathematical Function}

Define \(P^\impl:\N^4\to[0,1]\) recursively as follows:
\[
    P^\impl(s)=1\quad\text{if }s\in\Success,
\]
\[
    P^\impl(s)=0\quad\text{if }s\in\Failure,
\]
and, for \(s\in\Active\),
\[
    P^\impl(s)=
    \max_{a\in\Actions(s)}
    \sum_{s'}K(s'\mid s,a)P^\impl(s').
\]

This is the exact mathematical version of the recursive algorithm.  The implementation may store decimal constants as floating-point values, but the proof concerns the ideal exact recurrence with probabilities \(1/5\), \(1/2\), and \(3/10\).  The planned verifier script can compute these values exactly using rational arithmetic.

\begin{theorem}[Recursive correctness]
For every state \(s\in\N^4\),
\[
    P^\impl(s)=V^\opt(s).
\]
That is, the recursive value returned by the algorithm equals the optimal probability of eventually obtaining at least one tier-4 horse.
\end{theorem}

\begin{proof}
We prove the result by strong induction on the rank
\[
    \rank(s)=t_1+t_2+t_3+t_4.
\]

\textbf{Base cases.}
If \(s\in\Success\), then by definition
\[
    P^\impl(s)=1=V^\opt(s).
\]
If \(s\in\Failure\), then by definition
\[
    P^\impl(s)=0=V^\opt(s).
\]

These handle all terminal states.  They also include all states of sufficiently small rank for which no legal move exists.

\textbf{Inductive hypothesis.}
Assume that for some \(n\ge0\), the equality
\[
    P^\impl(r)=V^\opt(r)
\]
holds for every state \(r\) with
\[
    \rank(r)<n.
\]

\textbf{Inductive step.}
Let \(s\in\Active\) with \(\rank(s)=n\).  For every legal action \(a\in\Actions(s)\) and every successor state \(s'\) with \(K(s'\mid s,a)>0\), the rank-decrease lemma gives
\[
    \rank(s')=\rank(s)-1=n-1<n.
\]
Therefore, by the inductive hypothesis,
\[
    P^\impl(s')=V^\opt(s')
\]
for every possible successor \(s'\).

Using the recursive definition of \(P^\impl\),
\[
\begin{aligned}
    P^\impl(s)
    &=
    \max_{a\in\Actions(s)}
    \sum_{s'}K(s'\mid s,a)P^\impl(s')\\
    &=
    \max_{a\in\Actions(s)}
    \sum_{s'}K(s'\mid s,a)V^\opt(s')\\
    &=
    V^\opt(s),
\end{aligned}
\]
where the final equality is the Bellman optimality equation.

Thus \(P^\impl(s)=V^\opt(s)\) for every active state of rank \(n\).  By strong induction, the equality holds for all finite states \(s\in\N^4\).
\end{proof}

\section{Initial Tier-1 Interface}

The main interface begins with \(X\) tier-1 horses and no other horses:
\[
    s_X=(X,0,0,0).
\]

Define
\[
    f(X)=P^\impl(s_X)=V^\opt(s_X).
\]

Thus \(f(X)\) is the optimal probability of eventually obtaining at least one tier-4 horse from \(X\) tier-1 horses.

\section{Minimal Resource Lemma}

The code records the practical observation that fewer than eight tier-1 horses cannot produce a tier-4 horse, while exactly eight tier-1 horses can succeed only along the perfect upgrade path.  We now prove this formally.

\begin{definition}[Tier-weight potential]
Define the tier-weight potential
\[
    \weight(s)=t_1+2t_2+4t_3+8t_4.
\]
This counts the number of tier-1 units represented by the current state if each tier is assigned its binary merge value.
\end{definition}

\begin{lemma}[Tier-weight is nonincreasing]
For every legal transition \(s\to s'\),
\[
    \weight(s')\le \weight(s).
\]
Moreover, equality holds exactly for upgrade outcomes.
\end{lemma}

\begin{proof}
We check each action.

For \(a_1\), the upgrade transition is
\[
    2T_1\to 1T_2,
\]
which changes tier-weight from \(2\) to \(2\).  The non-upgrade tier-1 outcome changes tier-weight from \(2\) to \(1\).

For \(a_2\), the upgrade transition is
\[
    2T_2\to 1T_3,
\]
which changes tier-weight from \(4\) to \(4\).  The same-tier outcome changes tier-weight from \(4\) to \(2\), and the downgrade outcome changes tier-weight from \(4\) to \(1\).

For \(a_3\), the upgrade transition is
\[
    2T_3\to 1T_4,
\]
which changes tier-weight from \(8\) to \(8\).  The same-tier outcome changes tier-weight from \(8\) to \(4\), and the downgrade outcome changes tier-weight from \(8\) to \(2\).

Therefore \(\weight\) never increases, and equality occurs exactly on upgrade outcomes.
\end{proof}

\begin{lemma}[Fewer than eight tier-1 horses cannot succeed]
If \(X<8\), then
\[
    f(X)=0.
\]
\end{lemma}

\begin{proof}
Starting from \(s_X=(X,0,0,0)\),
\[
    \weight(s_X)=X.
\]
If \(X<8\), then the initial tier-weight is less than \(8\).  Since tier-weight is nonincreasing along every transition, no reachable state can have tier-weight at least \(8\).  But every success state has \(t_4\ge1\), hence has
\[
    \weight(s)\ge 8.
\]
Therefore no success state is reachable, and the optimal success probability is \(0\).
\end{proof}

\begin{lemma}[Exactly eight tier-1 horses succeed only by perfect upgrades]
Starting from \(s_8=(8,0,0,0)\), every successful trajectory consists entirely of upgrade outcomes.  Consequently,
\[
    f(8)=\left(\frac15\right)^7.
\]
\end{lemma}

\begin{proof}
The initial tier-weight is
\[
    \weight(s_8)=8.
\]
Any success state must have \(t_4\ge1\), hence tier-weight at least \(8\).  Since tier-weight is nonincreasing, any successful trajectory from \(s_8\) must preserve tier-weight at every step.  By the previous lemma, tier-weight is preserved exactly when every outcome is an upgrade.

It remains to count the number of upgrade outcomes required.  Starting from eight tier-1 horses, reaching one tier-4 horse by binary merging requires
\[
    8T_1 \to 4T_2 \to 2T_3 \to 1T_4.
\]
This consists of
\[
    4+2+1=7
\]
upgrade outcomes.  Each upgrade has probability \(1/5\), and these outcome probabilities multiply along the path.  Therefore
\[
    f(8)=\left(\frac15\right)^7.
\]
\end{proof}

\section{Empirical Scaling: What Is Proven and What Is Observed}

The implementation includes a log-log regression step intended to describe how \(f(X)\) scales over sampled values of \(X\).  This is useful empirical analysis, but it should be separated from the exact correctness theorem above.

\begin{definition}[Finite-window scaling fit]
Let \(I\subset\N\) be a finite set of sampled inputs with \(f(X)>0\) for all \(X\in I\).  A finite-window power-law fit estimates constants \(c>0\) and \(\alpha\in\R\) such that
\[
    f(X)\approx cX^\alpha
\]
for \(X\in I\), usually by linear regression on
\[
    \log f(X)\approx \log c+\alpha\log X.
\]
\end{definition}

\begin{proposition}[No positive power law can hold globally]
There do not exist constants \(c>0\), \(\alpha>0\), and \(X_0\in\N\) such that
\[
    f(X)=cX^\alpha
\]
for all \(X\ge X_0\).
\end{proposition}

\begin{proof}
For every \(X\), \(f(X)\) is a probability, so
\[
    0\le f(X)\le 1.
\]
If \(f(X)=cX^\alpha\) for all \(X\ge X_0\), with \(c>0\) and \(\alpha>0\), then
\[
    cX^\alpha\to\infty
\]
as \(X\to\infty\).  In particular, for sufficiently large \(X\), one has \(cX^\alpha>1\), contradicting \(f(X)\le1\).  Therefore no such global positive power law exists.
\end{proof}

\begin{remark}[Safe scaling language]
A statement such as
\[
    f(X)\propto X^{1.8}
\]
should be interpreted as a finite-range empirical fit unless a bounded saturating model is supplied.  A safer claim is:

\begin{quote}
Over the sampled finite range, the log-log regression of \(f(X)\) against \(X\) exhibits an approximate power-law exponent near \(\alpha\).
\end{quote}

This preserves the useful empirical observation without overstating it as a global theorem.
\end{remark}

\section{Proof-Tool Selection Map}

This project is also useful as a worked example of proof-tool selection.  The theorem's shape and mathematical material naturally select the proof tools.

\subsection{Claim shape}

The main claim is an algorithm-correctness theorem:
\[
    \forall s\in\N^4,\quad P^\impl(s)=V^\opt(s).
\]
This activates the following proof-tool nodes:

\[
    \text{Universal claim}
    \longrightarrow
    \text{arbitrary state}
    \longrightarrow
    \text{induction candidate}.
\]

\subsection{Mathematical material}

The object is a finite-horizon recursive decision process.  This activates:

\[
    \text{State-transition graph}
    \longrightarrow
    \text{rank / monovariant}
    \longrightarrow
    \text{well-founded recursion}.
\]

The decision component activates:

\[
    \text{choice of action}
    \longrightarrow
    \text{Bellman recurrence}
    \longrightarrow
    \text{optimal substructure}.
\]

\subsection{Selected proof tools}

The final selected tools are:

\begin{enumerate}
    \item \textbf{Rank monovariant:} \(\rank(s)\) decreases after every legal move.
    \item \textbf{Strong induction:} prove correctness over increasing total horse count.
    \item \textbf{Bellman optimality:} first action plus optimal continuation.
    \item \textbf{Potential function:} \(\weight(s)\) proves the eight-horse resource bound.
    \item \textbf{Boundedness audit:} probability bounded by \(1\) prevents global positive power-law scaling.
\end{enumerate}

\section{Implementation Obligations for the Verification Script}

The next computational artifact should verify the proof numerically and exactly.  A recommended Python verifier should include:

\begin{enumerate}
    \item Exact rational constants:
    \[
        p_+=\frac15,\qquad p_0=\frac12,\qquad p_-=\frac3{10}.
    \]
    \item A recursive memoized solver for \(P^\impl(s)\).
    \item A bottom-up dynamic-programming solver ordered by \(\rank(s)\).
    \item Equality checks between recursive and bottom-up values for all states up to a chosen rank bound.
    \item The exact check
    \[
        f(8)=\left(\frac15\right)^7.
    \]
    \item A Monte Carlo simulator used only as a sanity check, not as proof.
    \item A finite-window scaling report that labels fitted exponents as empirical.
\end{enumerate}

\section{Conclusion}

The recursive horse-combination calculator is correct because the process is a finite acyclic Markov decision process under the rank
\[
    \rank(s)=t_1+t_2+t_3+t_4.
\]
Every legal action reduces this rank by one, so the recursive value function is well-founded.  At each active state, the optimal strategy chooses the legal action with maximal expected optimal continuation value.  This is exactly the Bellman recurrence implemented by the algorithm.

The exact theorem is therefore:

\[
    \boxed{
    P^\impl(s)=V^\opt(s)
    \quad
    \text{for every finite horse state }s.
    }
\]

The minimal-resource theorem adds:

\[
    \boxed{
    f(X)=0\text{ for }X<8,
    \qquad
    f(8)=\left(\frac15\right)^7.
    }
\]

Empirical scaling analysis remains valuable, but it belongs to the computational-observation layer rather than the formal theorem layer.

\end{document}
